% GENERATED FILE. Edit canonical lessons, then run npm run generate:books.
% canonical-source-hash: fnv1a64:2e462c5e01ce821e
% canonical-lessons: SA-C400-parikramati, SA-C400-japamala, SA-C400-rudraksah, SA-C400-rathayatra, SA-C400-aharati

\chapter{Round the Shrine}
\label{ch:sa-round-the-shrine}
\hlchaptermodality{400}

\begin{chapteropening}
\textbf{By the end of this chapter:} \emph{I can say how pilgrims go round a shrine and what they bring.}

Complete the last lesson of chapter 400: I can say how pilgrims go round a shrine and what they bring.
\end{chapteropening}

\section[\texorpdfstring{\sk{परिक्रामति} (parikrāmati)}{परिक्रामति (parikrāmati)}]{\sk{परिक्रामति} (parikrāmati) — walks around; goes round}

\label{lesson:SA-C400-parikramati}



\begin{quote}
Before the new one: say the Sanskrit for an arched gateway, then the Sanskrit for a pennant.
\end{quote}

\subsection*{You'll want to know: \sk{परिक्रामति}}
\textbf{\sk{परिक्रामति}} — \emph{parikrāmati} — \textquotedblleft{}walks around; goes round\textquotedblright{}.

\begin{quote}
\sk{यात्रिकाः} \sk{मन्दिरं} \sk{परिक्रामन्ति।}
\end{quote}

\begin{quote}
\emph{yātrikāḥ mandiraṃ parikrāmanti.}
\end{quote}

\begin{quote}
The pilgrims walk around the temple.
\end{quote}

Past (he …): \sk{सः} \sk{मन्दिरं} \sk{परिक्रान्तवान्}, \emph{saḥ mandiraṃ parikrāntavān}.

\begin{cousinweb}
\textbf{\sk{परिक्रामति}} (\emph{parikrāmati}) is built from \emph{pari}, \textquotedblleft{}around\textquotedblright{}, and the root √kram, \textquotedblleft{}to step, to stride\textquotedblright{}; MW gives \emph{pari-kram} as to walk or step round, with the participle \emph{parikrānta} (MW).
\end{cousinweb}

\begin{grammarlens}[title={What you've built}]
One more word for getting around.
\end{grammarlens}

\subsection*{Your turn}
\begin{itemize}
\raggedright
  \item \emph{Say it:} \emph{parikrāmati}
  \item \emph{Say it:} \emph{parikrāmati}, once more
  \item \emph{Say it:} \emph{yātrikāḥ mandiraṃ parikrāmanti.} — The pilgrims walk around the temple.
  \item \emph{Recall:} say the Sanskrit for an arched gateway, then the Sanskrit for a pennant, then say \emph{parikrāmati} again
\end{itemize}

\begin{tcolorbox}[breakable,colback=teal!4,colframe=teal!35!black,title={Before you move on}]
What does \sk{परिक्रामति} mean? (\textquotedblleft{}Walks around; goes round\textquotedblright{}.) Say it once more.
\end{tcolorbox}
\section[\texorpdfstring{\sk{जपमाला} (japamālā)}{जपमाला (japamālā)}]{\sk{जपमाला} (japamālā) — a string of prayer beads}

\label{lesson:SA-C400-japamala}



\begin{quote}
Before the new one: say the Sanskrit for a pennant, then the Sanskrit for walks around.
\end{quote}

\subsection*{You'll want to know: \sk{जपमाला}}
\textbf{\sk{जपमाला}} — \emph{japamālā} — \textquotedblleft{}a string of prayer beads\textquotedblright{}.

\begin{quote}
\sk{वृद्धः} \sk{हस्ते} \sk{जपमालां} \sk{धारयति।}
\end{quote}

\begin{quote}
\emph{vṛddhaḥ haste japamālāṃ dhārayati.}
\end{quote}

\begin{quote}
The old man holds prayer beads in his hand.
\end{quote}

A feminine noun.

\begin{cousinweb}
\textbf{\sk{जपमाला}} (\emph{japamālā}) is built from \emph{japa}, \textquotedblleft{}the quiet repeating of prayers\textquotedblright{}, from the root √jap, \textquotedblleft{}to mutter\textquotedblright{}, and \emph{mālā}, \textquotedblleft{}a garland\textquotedblright{}; MW gives the compound as a rosary for counting prayers (MW).
\end{cousinweb}

\begin{grammarlens}[title={What you've built}]
One more word for getting around.
\end{grammarlens}

\subsection*{Your turn}
\begin{itemize}
\raggedright
  \item \emph{Say it:} \emph{japamālā}
  \item \emph{Say it:} \emph{japamālā}, once more
  \item \emph{Say it:} \emph{vṛddhaḥ haste japamālāṃ dhārayati.} — The old man holds prayer beads in his hand.
  \item \emph{Recall:} say the Sanskrit for a pennant, then the Sanskrit for walks around, then say \emph{japamālā} again
\end{itemize}

\begin{tcolorbox}[breakable,colback=teal!4,colframe=teal!35!black,title={Before you move on}]
What does \sk{जपमाला} mean? (\textquotedblleft{}A string of prayer beads\textquotedblright{}.) Say it once more.
\end{tcolorbox}
\section[\texorpdfstring{\sk{रुद्राक्षः} (rudrākṣaḥ)}{रुद्राक्षः (rudrākṣaḥ)}]{\sk{रुद्राक्षः} (rudrākṣaḥ) — a rudraksha bead}

\label{lesson:SA-C400-rudraksah}



\begin{quote}
Before the new one: say the Sanskrit for walks around, then the Sanskrit for prayer beads.
\end{quote}

\subsection*{You'll want to know: \sk{रुद्राक्षः}}
\textbf{\sk{रुद्राक्षः}} — \emph{rudrākṣaḥ} — \textquotedblleft{}a rudraksha bead\textquotedblright{}.

\begin{quote}
\sk{जपमालायां} \sk{रुद्राक्षाः} \sk{सन्ति।}
\end{quote}

\begin{quote}
\emph{japamālāyāṃ rudrākṣāḥ santi.}
\end{quote}

\begin{quote}
There are rudraksha beads on the string of prayer beads.
\end{quote}

A masculine noun.

\begin{cousinweb}
\textbf{\sk{रुद्राक्षः}} (\emph{rudrākṣaḥ}) is built from \emph{Rudra}, a name of Śiva, and \emph{akṣa}, \textquotedblleft{}an eye\textquotedblright{}: MW gives it as \textquotedblleft{}Rudra-eyed\textquotedblright{}, a tree whose berries are strung into rosaries (MW).
\end{cousinweb}

\begin{grammarlens}[title={What you've built}]
One more word for getting around.
\end{grammarlens}

\subsection*{Your turn}
\begin{itemize}
\raggedright
  \item \emph{Say it:} \emph{rudrākṣaḥ}
  \item \emph{Say it:} \emph{rudrākṣaḥ}, once more
  \item \emph{Say it:} \emph{japamālāyāṃ rudrākṣāḥ santi.} — There are rudraksha beads on the string of prayer beads.
  \item \emph{Recall:} say the Sanskrit for walks around, then the Sanskrit for prayer beads, then say \emph{rudrākṣaḥ} again
\end{itemize}

\begin{tcolorbox}[breakable,colback=teal!4,colframe=teal!35!black,title={Before you move on}]
What does \sk{रुद्राक्षः} mean? (\textquotedblleft{}A rudraksha bead\textquotedblright{}.) Say it once more.
\end{tcolorbox}
\section[\texorpdfstring{\sk{रथयात्रा} (rathayātrā)}{रथयात्रा (rathayātrā)}]{\sk{रथयात्रा} (rathayātrā) — a chariot procession}

\label{lesson:SA-C400-rathayatra}



\begin{quote}
Before the new one: say the Sanskrit for prayer beads, then the Sanskrit for a rudraksha bead.
\end{quote}

\subsection*{You'll want to know: \sk{रथयात्रा}}
\textbf{\sk{रथयात्रा}} — \emph{rathayātrā} — \textquotedblleft{}a chariot procession\textquotedblright{}.

\begin{quote}
\sk{रथयात्रायां} \sk{जनाः} \sk{रथं} \sk{कर्षन्ति।}
\end{quote}

\begin{quote}
\emph{rathayātrāyāṃ janāḥ rathaṃ karṣanti.}
\end{quote}

\begin{quote}
In the chariot procession people pull the chariot.
\end{quote}

A feminine noun.

\begin{cousinweb}
\textbf{\sk{रथयात्रा}} (\emph{rathayātrā}) is built from \emph{ratha}, \textquotedblleft{}a chariot\textquotedblright{}, and \emph{yātrā}, \textquotedblleft{}a journey, a procession\textquotedblright{}; MW gives the compound as the festive procession of an image on a car (MW).
\end{cousinweb}

\begin{grammarlens}[title={What you've built}]
One more word for getting around.
\end{grammarlens}

\subsection*{Your turn}
\begin{itemize}
\raggedright
  \item \emph{Say it:} \emph{rathayātrā}
  \item \emph{Say it:} \emph{rathayātrā}, once more
  \item \emph{Say it:} \emph{rathayātrāyāṃ janāḥ rathaṃ karṣanti.} — In the chariot procession people pull the chariot.
  \item \emph{Recall:} say the Sanskrit for prayer beads, then the Sanskrit for a rudraksha bead, then say \emph{rathayātrā} again
\end{itemize}

\begin{tcolorbox}[breakable,colback=teal!4,colframe=teal!35!black,title={Before you move on}]
What does \sk{रथयात्रा} mean? (\textquotedblleft{}A chariot procession\textquotedblright{}.) Say it once more.
\end{tcolorbox}
\section[\texorpdfstring{\sk{आहरति} (āharati)}{आहरति (āharati)}]{\sk{आहरति} (āharati) — fetches; brings}

\label{lesson:SA-C400-aharati}



\begin{quote}
Before the new one: say the Sanskrit for a rudraksha bead, then the Sanskrit for a chariot procession.
\end{quote}

\subsection*{You'll want to know: \sk{आहरति}}
\textbf{\sk{आहरति}} — \emph{āharati} — \textquotedblleft{}fetches; brings\textquotedblright{}.

\begin{quote}
\sk{बालिका} \sk{उद्यानात्} \sk{पुष्पाणि} \sk{आहरति।}
\end{quote}

\begin{quote}
\emph{bālikā udyānāt puṣpāṇi āharati.}
\end{quote}

\begin{quote}
The girl fetches flowers from the garden.
\end{quote}

Past (she …): \sk{सा} \sk{पुष्पाणि} \sk{आहृतवती}, \emph{sā puṣpāṇi āhṛtavatī}.

\begin{cousinweb}
\textbf{\sk{आहरति}} (\emph{āharati}) is built from \emph{ā}, \textquotedblleft{}towards\textquotedblright{}, and the root √hṛ, \textquotedblleft{}to carry\textquotedblright{}, which also gives \emph{harati}, \textquotedblleft{}takes away\textquotedblright{}; MW gives \emph{ā-hṛ} as to fetch or bring, with the participle \emph{āhṛta} (MW).
\end{cousinweb}

\begin{grammarlens}[title={What you've built}]
One more word for getting around.
\end{grammarlens}

\subsection*{Your turn}
\begin{itemize}
\raggedright
  \item \emph{Say it:} \emph{āharati}
  \item \emph{Say it:} \emph{āharati}, once more
  \item \emph{Say it:} \emph{bālikā udyānāt puṣpāṇi āharati.} — The girl fetches flowers from the garden.
  \item \emph{Recall:} say the Sanskrit for a rudraksha bead, then the Sanskrit for a chariot procession, then say \emph{āharati} again
\end{itemize}

\begin{tcolorbox}[breakable,colback=teal!4,colframe=teal!35!black,title={Before you move on}]
What does \sk{आहरति} mean? (\textquotedblleft{}Fetches; brings\textquotedblright{}.) Say it once more.
\end{tcolorbox}
